\begin{center}
    \textbf{Vilnius University} \\
    \textbf{Systems Biology} \\
    
    \vspace{0.5cm}
    
    \textbf{ENHANCING SINGLE-CELL RNA-SEQ ANALYSIS THROUGH THE INTEGRATION OF UNACCOUNTED INTERGENIC REGIONS} \\
    
    \vspace{1cm}
    
    \textbf{Juozapas Ivanauskas} \\
    
    \vspace{1cm}
    
    \textbf{Abstract} \\
    
\end{center}

Single cell RNA sequencing data contain reads that map to the genome but are not assigned to any gene.
The aim of this project was to investigate such reads, to reveal their origins and if they are not technical artifacts,
to include them into downstream analysis.
In this project unassigned reads from fourteen datasets were analyzed.
It was shown, that majority of those comes from antisense regions of known genes, and hypothesis was suggested,
that those are artifacts produced during reverse transcription reaction.
Some of the reads were mapped to intergenic regions without any genes on any strand,
and those are likely originating from unannotated (novel) genes.
Finally, a part of reads were mapped to gene regions but still left unassigned.
To include those into downstream analysis, transcriptomic reference was enhanced, by resolving some overlapping genes,
extending 3' ends of part of the genes and adding some genes from different transcriptomic annotations.
Such approach allowed to have on average 1.2\% more counts in the downstream analysis.
